\section{A sharp weighted transfer from bounded projective masks}
\label{sec:tpc49-weighted-transfer}

Let \(I,J\) be finite sets, let \(K:I\times J\to\C\), and let
\(a=(a_i)_{i\in I}\), \(b=(b_j)_{j\in J}\) be nonnegative weights.
Zero weights may be deleted.  Write
\begin{equation}
 A:=D_aKD_b,
 \qquad A_{ij}=a_iK_{ij}b_j.
 \label{eq:tpc49-weighted-matrix}
\end{equation}
Complex unary row and column factors cause no loss of generality here:
their moduli are placed in \(a,b\), while their unit phases are absorbed
into \(K\).  Such unit-modulus row and column Schur factors preserve every
bounded-projective cost used below.
When \(A\ne0\), define its weighted survival parameter by
\begin{equation}
 \delta_w(K;a,b)
 :=\frac{\|A\|_{S_2}^2}{\|a\|_2^2\|b\|_2^2}
 =\frac{\sum_{i,j}a_i^2|K_{ij}|^2b_j^2}
 {\left(\sum_i a_i^2\right)\left(\sum_jb_j^2\right)}.
 \label{eq:tpc49-weighted-admission}
\end{equation}
For a \(0\)-\(1\) mask it lies in \([0,1]\) and is literally the
proportion of product \(\ell^2\)-mass lying on admitted pairs.

Fix an exact bounded-projective representation
\begin{equation}
 \mathscr R:\qquad
 K_{ij}=\int_\Omega p_\omega(i)q_\omega(j)\,d\nu(\omega),
 \label{eq:tpc49-bounded-representation}
\end{equation}
and put
\begin{align}
 P_{\mathscr R}(K)
 &:=\int_\Omega\|p_\omega\|_\infty
       \|q_\omega\|_\infty\,d|\nu|(\omega),
 \label{eq:tpc49-fixed-bounded-cost}\\
 \cE_{\mathscr R}(K;a,b)
 &:=\int_\Omega
 \|D_ap_\omega\|_2\|D_b\overline{q_\omega}\|_2
 \,d|\nu|(\omega).
 \label{eq:tpc49-weighted-Hilbert-envelope}
\end{align}
The conjugation merely matches the operator convention and has no
effect on any norm.

\begin{theorem}[Weighted projective-to-atomic transfer]
\label{thm:tpc49-weighted-transfer}
For every representation \eqref{eq:tpc49-bounded-representation},
\begin{align}
 \|A\|_{S_1}
 &\le \cE_{\mathscr R}(K;a,b)
 \le P_{\mathscr R}(K)\|a\|_2\|b\|_2,
 \label{eq:tpc49-fixed-weighted-transfer}\\
 \frac{\cE_{\mathscr R}(K;a,b)^2}{\|A\|_{S_2}^2}
 &\le\frac{P_{\mathscr R}(K)^2}{\delta_w(K;a,b)}.
 \label{eq:tpc49-fixed-weighted-energy}
\end{align}
Taking the infimum in the bounded projective cost gives
\begin{equation}
 \boxed{
 r_{\rm nuc}(A)
 \le\min\!\left\{\operatorname{rank}A,
 \frac{\|K\|_{\pi,\infty}^2}{\delta_w(K;a,b)}
 \right\}.}
 \label{eq:tpc49-optimal-weighted-transfer}
\end{equation}
If \(a_i,b_j>0\) on the active sets and one takes the infimum of
\eqref{eq:tpc49-weighted-Hilbert-envelope} over all exact
representations there, the result is exactly \(\|A\|_{S_1}\).
This unrestricted infimum is not asserted to coincide with the infimum
over the narrower inherited admissible TPC--48 representation class.
\end{theorem}

\begin{proof}
After multiplying \eqref{eq:tpc49-bounded-representation} by the two
diagonal weights,
\begin{equation}
 A=\int_\Omega
 (D_ap_\omega)\otimes
 \overline{(D_b\overline{q_\omega})}\,d\nu(\omega).
 \label{eq:tpc49-weighted-rank-one-integral}
\end{equation}
The nuclear norm is at most the integral of the rank-one nuclear norms,
which is the first inequality in
\eqref{eq:tpc49-fixed-weighted-transfer}.  Pointwise,
\begin{equation}
 \|D_ap_\omega\|_2
 \|D_b\overline{q_\omega}\|_2
 \le\|a\|_2\|b\|_2
 \|p_\omega\|_\infty\|q_\omega\|_\infty.
 \label{eq:tpc49-layer-sup-to-Hilbert}
\end{equation}
Integration proves the second inequality.  Divide its square by
 \(\|A\|_{S_2}^2=\delta_w\|a\|_2^2\|b\|_2^2\) to obtain
\eqref{eq:tpc49-fixed-weighted-energy}.  Infimize and combine with
\cref{thm:tpc49-Hilbert-projective} for
\eqref{eq:tpc49-optimal-weighted-transfer}.  Finally, when the weights
are positive, every rank-one decomposition of \(A\) can be divided by
\(a_i,b_j\) to give a representation of \(K\) on the active rectangle;
the exact infimum is therefore the nuclear norm.
\end{proof}

\begin{remark}[Sufficient and exact criteria]
\label{rem:tpc49-density-not-necessary}
The exact intrinsic criterion is \(r_{\rm nuc}(A)\), and for a fixed
representation it is
\(c_{\mathscr R}(A)^2r_{\rm nuc}(A)\).
The quantity \(\|K\|_{\pi,\infty}^2/\delta_w\) is a computable
sufficient upper bound.  A lower bound on \(\delta_w\) is not necessary
for a particular matrix: a sparse rank-one rectangle can have tiny
density and \(r_{\rm nuc}=1\).  The density dependence is instead sharp
in the minimax class controlled only by bounded projective cost and
\(\delta_w\), as \cref{prop:tpc49-permutation-sharp} shows.
\end{remark}

\subsection{A uniform degree criterion}

For a \(0\)-\(1\) mask \(K\) with \(|J|=N\), suppose every row has at
most \(\epsilon N\) forbidden entries.  Define the flatness of the
right weight by
\begin{equation}
 \operatorname{Flat}(b)
 :=\frac{N\|b\|_\infty^2}{\|b\|_2^2}\ge1.
 \label{eq:tpc49-weight-flatness}
\end{equation}

\begin{corollary}[Degree plus flatness]
\label{cor:tpc49-degree-flatness}
For arbitrary nonzero left weights \(a\),
\begin{equation}
 \delta_w(K;a,b)
 \ge1-\epsilon\operatorname{Flat}(b).
 \label{eq:tpc49-degree-density}
\end{equation}
Consequently, if the right side is positive,
\begin{equation}
 r_{\rm nuc}(D_aKD_b)
 \le\frac{\|K\|_{\pi,\infty}^2}
 {1-\epsilon\operatorname{Flat}(b)}.
 \label{eq:tpc49-degree-effective-rank}
\end{equation}
The same statement holds with the two sides exchanged.
\end{corollary}

\begin{proof}
For each row \(i\), its forbidden \(b^2\)-mass is at most
\(\epsilon N\|b\|_\infty^2
=\epsilon\operatorname{Flat}(b)\|b\|_2^2\).
Average the complementary inequality with weights \(a_i^2\), then use
\cref{thm:tpc49-weighted-transfer}.
\end{proof}

This criterion tolerates arbitrary concentration in the source weight
provided the opposite-row weight is flat.  For the TPC joint multiplier,
flatness and admitted density of the active smooth/residue cell are new
conditions: the inherited interface gives upper seminorms and exact
identities, but no lower-density theorem.

\subsection{The sharp permutation boundary}

\begin{proposition}[Permutation equality]
\label{prop:tpc49-permutation-sharp}
Let \(P_n\) be any \(n\times n\) permutation matrix and take
\(a=b=\one\).  Then
\begin{equation}
 \|P_n\|_{\pi,\infty}=1,
 \qquad
 \delta_w(P_n;\one,\one)=\frac1n,
 \qquad
 \|P_n\|_{S_1}=n,
 \qquad
 \|P_n\|_{S_2}^2=n.
 \label{eq:tpc49-permutation-data}
\end{equation}
Thus \(r_{\rm nuc}(P_n)=n\), and
\eqref{eq:tpc49-optimal-weighted-transfer} is an equality.
\end{proposition}

\begin{proof}
Permutations preserve singular values, so all \(n\) singular values are
one.  For the identity matrix,
\begin{equation}
 I_n=2^{-n}\sum_{\varepsilon\in\{\pm1\}^n}
 \varepsilon\varepsilon^{\mathsf T}.
 \label{eq:tpc49-Rademacher-identity}
\end{equation}
Every factor has supremum norm one, giving bounded projective cost at
most one; the maximum entry gives the reverse inequality.  Permuting a
factor preserves the cost.  The remaining formulas follow by counting
the \(n\) nonzero entries.
\end{proof}

The example is nonnegative, one-sign, and has an orthogonal singular
basis.  None of those qualitative properties replaces a quantitative
effective-rank or density estimate.
